\clearpage
\section{06 SQL 查询：中文学习笔记}\label{06_sql_queries_zh}

CSC5041 · 第一次课堂测验复习 · 2026-10-19

\subsection{阅读说明与学习目标}\label{06_sql_queries_zh-ux9605ux8bfbux8bf4ux660eux4e0eux5b66ux4e60ux76eeux6807}

本篇接续 \hyperref[05_relational_algebra_zh]{05 关系代数}，主要依据教师 \href{../../../slides/SQLpart1.pptx}{SQLpart1.pptx} 第 3--58 页，以及 \href{../../../slides/SQLpart2.pptx}{SQLpart2.pptx} 第 2--24 页中与查询相关的内容。下文用 P1/P2 区分两份课件，页码按幻灯片顺序计数。

教学计划第 1 页 Week 5 包含 SELECT--FROM--WHERE、排序、聚合、分组与嵌套查询；第 2 页 Tutorial 2 明列 GROUP BY、HAVING、NOT EXISTS、NOT IN with NULLs。因此不能只读文件名为 Part 1 的课件就停止。高级量化比较、DDL/DML、视图及 PL/pgSQL 的边界会另标，不把整份 Part 2 都认定为 Test 1 必考内容。

学完后，应能选择正确的结果粒度，写完整连接条件，处理重复值、NULL、分组、存在性和反存在性。所有例子使用自编数据，\textbf{查询结果是静态推演，未在 PostgreSQL 或其他数据库中执行}。补充官方文档用于核对语义，不决定教师考试深度。

\Needspace{11cm}
\subsection{1. 六个常用子句：书写顺序与逻辑顺序}\label{06_sql_queries_zh-ux516dux4e2aux5e38ux7528ux5b50ux53e5ux4e66ux5199ux987aux5e8fux4e0eux903bux8f91ux987aux5e8f}

\StudyVisual{sqlflow}{zh}

来源：P1 第 13、26--27、44--45 页；P2 第 18、22--24 页。

\begin{StudyCode}
SELECT 输出列或表达式
FROM 输入表及连接
WHERE 行条件
GROUP BY 分组属性
HAVING 组条件
ORDER BY 排序属性;
\end{StudyCode}

WHERE、GROUP BY、HAVING、ORDER BY 按需要使用。P2 第 24 页展示的是本课查询框架，不是现代 SQL 所有语法的完整上限。

理解常规查询时可按以下\textbf{逻辑阶段}阅读：

\begin{StudyCode}
FROM / JOIN → WHERE → GROUP BY / 聚合 → HAVING
            → SELECT → DISTINCT（若指定）→ ORDER BY
\end{StudyCode}

这不是优化器必须采用的物理执行顺序。WHERE 在组形成前筛行，不能在同一查询层直接用 \texttt{AVG(\allowbreak{}score)\allowbreak{}} 当行条件；输出别名也不要直接拿到同层 WHERE 中使用。ORDER BY 则可使用输出别名。

SQL 查询通常是声明式的：说明要什么结果。SELECT 中写一个计算表达式，只改变输出，不等于 UPDATE 修改原始数据。

\subsection{2. 贯穿例子：在 05 的基础上加入学期与缺失成绩}\label{06_sql_queries_zh-ux8d2fux7a7fux4f8bux5b50ux5728-05-ux7684ux57faux7840ux4e0aux52a0ux5165ux5b66ux671fux4e0eux7f3aux5931ux6210ux7ee9}

以下为自编输入；\textbf{不是 Tutorial 原模式}，也不沿用 05 的 Takes 作为完整选课表。

\begin{StudyCode}
Student(sid, sname, major)                 PK{sid}
Course(cid, title, dept, credits)          PK{cid}
Enrolment(sid, cid, term, score)           PK{sid,cid,term}
Enrolment.sid → Student.sid; Enrolment.cid → Course.cid
\end{StudyCode}

除 score 外各列均必填；score 是可空数值，NULL 表示成绩尚未获得，不是 0。term 是字符串。

Student：

\begin{longtable}[]{@{}lll@{}}
\toprule\noalign{}
sid & sname & major \\
\midrule\noalign{}
\endhead
\bottomrule\noalign{}
\endlastfoot
S1 & Amy & CS \\
S2 & Ben & Math \\
S3 & Amy & CS \\
S4 & Dana & Math \\
\end{longtable}

Course：

\begin{longtable}[]{@{}llll@{}}
\toprule\noalign{}
cid & title & dept & credits \\
\midrule\noalign{}
\endhead
\bottomrule\noalign{}
\endlastfoot
C1 & Databases & CS & 4 \\
C2 & AI & CS & 3 \\
C3 & Linear Algebra & Math & 4 \\
\end{longtable}

Enrolment：

\begin{longtable}[]{@{}llll@{}}
\toprule\noalign{}
sid & cid & term & score \\
\midrule\noalign{}
\endhead
\bottomrule\noalign{}
\endlastfoot
S1 & C1 & 2026F & 90 \\
S1 & C2 & 2026F & 80 \\
S1 & C1 & 2027S & NULL \\
S2 & C1 & 2026F & NULL \\
S2 & C3 & 2026F & 80 \\
S3 & C3 & 2026F & NULL \\
\end{longtable}

共有 6 次选课、3 个非空成绩、2 种不同非空成绩（80、90）。S1 重修 C1；S1 与 S3 同名；S4 无选课。后续除明确的假设变更外，都独立使用这份输入。

\subsection{3. 基本查询、别名、字符串与 DISTINCT}\label{06_sql_queries_zh-ux57faux672cux67e5ux8be2ux522bux540dux5b57ux7b26ux4e32ux4e0e-distinct}

来源：P1 第 12--13、18--25、28--40 页。

\Needspace{75pt}
\begin{lstlisting}[language=SQL]
SELECT c.cid, c.title, c.credits * 2 AS doubled_credits
FROM Course AS c
WHERE c.dept = 'CS' AND c.credits > 3;
\end{lstlisting}

结果为 \texttt{(\allowbreak{}C1,\allowbreak{}Databases,\allowbreak{}8)\allowbreak{}}，Course 中的 credits 仍是 4。字符串用半角单引号；字符串内的单引号加倍，例如 \texttt{\textquotesingle{}O\textquotesingle{}\textquotesingle{}Brien\textquotesingle{}}。不要把中文弯引号从 PPT 直接复制到代码中。

\begin{longtable}[]{@{}ll@{}}
\toprule\noalign{}
写法 & 含义 \\
\midrule\noalign{}
\endhead
\bottomrule\noalign{}
\endlastfoot
\texttt{=}、\texttt{\textless{}\textgreater{}}、\texttt{\textless{}}、\texttt{\textless{}=}、\texttt{\textgreater{}}、\texttt{\textgreater{}=} & 比较运算；本篇用 \texttt{\textless{}\textgreater{}} 表示不等 \\
\texttt{IN\ (\textquotesingle{}CS\textquotesingle{},\textquotesingle{}Math\textquotesingle{})} & 属于给定值列表 \\
\texttt{LIKE\ \textquotesingle{}A\%\textquotesingle{}} & 以 A 开头，\texttt{\%} 可匹配零个或多个字符 \\
\texttt{LIKE\ \textquotesingle{}\_e\%\textquotesingle{}} & 第二个字符是 e，\texttt{\_} 匹配一个字符 \\
\texttt{BETWEEN\ 3\ AND\ 4} & 含边界的范围比较 \\
\end{longtable}

BETWEEN 是补充写法。混用 AND、OR 时主动加括号，例如 \texttt{(major=\textquotesingle{}CS\textquotesingle{}\ OR\ major=\textquotesingle{}Math\textquotesingle{})\ AND\ sid\textless{}\textgreater{}\textquotesingle{}S4\textquotesingle{}}，避免把组合条件读错。

\subsubsection{SQL 默认保留输出重复}\label{06_sql_queries_zh-sql-ux9ed8ux8ba4ux4fddux7559ux8f93ux51faux91cdux590d}

\Needspace{60pt}
\begin{lstlisting}[language=SQL]
SELECT s.sname FROM Student AS s;
SELECT DISTINCT s.sname FROM Student AS s;
\end{lstlisting}

第一式输出 Amy 两次、Ben 一次、Dana 一次；第二式只有三种姓名。没有 ORDER BY 时不要假定这些行按某个顺序返回。\texttt{DISTINCT\ sid,\allowbreak{}sname} 对\textbf{整个输出属性组}去重，不是单独对每列去重。

只要输出姓名，不表示姓名是键；连接与集合判断仍应优先保留 sid。

\Needspace{11cm}
\subsection{4. NULL 与三值逻辑}\label{06_sql_queries_zh-null-ux4e0eux4e09ux503cux903bux8f91}

\StudyVisual{null}{zh}

来源：P1 第 18、35--36 页；补充核对 \href{https://www.postgresql.org/docs/18/functions-comparison.html}{PostgreSQL 比较运算文档}。

普通比较遇到 NULL 通常得到 \textbf{UNKNOWN}。WHERE 只保留条件为 TRUE 的行；FALSE 与 UNKNOWN 都不保留，但二者不是同一个逻辑值。

\begin{longtable}[]{@{}ll@{}}
\toprule\noalign{}
表达式 & 结果 \\
\midrule\noalign{}
\endhead
\bottomrule\noalign{}
\endlastfoot
\texttt{NULL\ =\ NULL}、\texttt{NULL\ \textless{}\textgreater{}\ NULL}、\texttt{NULL\ \textgreater{}\ 80} & UNKNOWN \\
\texttt{NOT\ UNKNOWN} & UNKNOWN \\
\texttt{TRUE\ AND\ UNKNOWN} & UNKNOWN \\
\texttt{FALSE\ AND\ UNKNOWN} & FALSE \\
\texttt{TRUE\ OR\ UNKNOWN} & TRUE \\
\texttt{FALSE\ OR\ UNKNOWN} & UNKNOWN \\
\end{longtable}

\textbf{课件措辞辨析}：P1 第 18 页的 ``Comparison against NULL yields FALSE'' 不应按字面记成逻辑定律。第 35 页明确给出 UNKNOWN；``被 WHERE 过滤掉''不等于``表达式本身为 FALSE''。

\Needspace{75pt}
\begin{lstlisting}[language=SQL]
SELECT sid, cid, term
FROM Enrolment
WHERE score IS NULL;
\end{lstlisting}

得到 \texttt{(\allowbreak{}S1,\allowbreak{}C1,\allowbreak{}2027S)\allowbreak{}}、\texttt{(\allowbreak{}S2,\allowbreak{}C1,\allowbreak{}2026F)\allowbreak{}}、\texttt{(\allowbreak{}S3,\allowbreak{}C3,\allowbreak{}2026F)\allowbreak{}}，结果无保证顺序。不能用 \texttt{score\ =\ NULL} 替代；也不能用 \texttt{score\ \textless{}\textgreater{}\ 90} 找出未出分记录，因为这些记录的比较结果是 UNKNOWN。

若要``不是 90 分，或者未出分''，明确写 \texttt{score\ \textless{}\textgreater{}\ 90\ OR\ score\ IS\ NULL}。数值 0 是一个已知成绩，不应无需求依据地把 NULL 改成 0。

\subsection{5. 多表连接：先确定完整配对条件}\label{06_sql_queries_zh-ux591aux8868ux8fdeux63a5ux5148ux786eux5b9aux5b8cux6574ux914dux5bf9ux6761ux4ef6}

来源：P1 第 41--49、52--58 页。

\Needspace{105pt}
\begin{lstlisting}[language=SQL]
SELECT s.sid, s.sname, c.title, e.term
FROM Student AS s
JOIN Enrolment AS e ON e.sid = s.sid
JOIN Course AS c ON c.cid = e.cid
WHERE c.dept = 'CS';
\end{lstlisting}

本例得到 4 行：S1 的 C1 两个学期、C2 一次，以及 S2 的 C1 一次。若只输出姓名，会有三行 Amy、一行 Ben；题目要求不同学生时，可输出 \texttt{DISTINCT\ s.\allowbreak{}sid,\allowbreak{}s.\allowbreak{}sname}，而不是掩盖错误连接。

遗漏连接条件可能产生笛卡尔积；两张有 name 的表写 \texttt{name=name} 还可能因歧义报错。用表别名限定列名，也让每条联系更容易检查。

\subsubsection{NATURAL JOIN 的风险}\label{06_sql_queries_zh-natural-join-ux7684ux98ceux9669}

它匹配\textbf{所有同名列}并合并它们。P1 第 55--57 页说明教师院系与课程院系未必相同，不能无意增加相等条件。优先用清楚的 JOIN\ldots ON 表达已确认的业务联系。

Tutorial 2 的 Enrolled 与 Teaches 都有 cid、term，教师与选课的关联要同时匹配二者：\texttt{e.\allowbreak{}cid=t.\allowbreak{}cid\ AND\ e.\allowbreak{}term=t.\allowbreak{}term}。只按 cid 会把不同学期混在一起。本篇没有 Teaching 表，因此不把该条件当作可直接对本例运行的完整查询。

对 INNER JOIN，某些条件可在 ON 与 WHERE 间等价调整；不要把它理解为固定的物理执行先后，也不要推广为所有外连接都等价，外连接会涉及补 NULL 的行。

\subsection{6. 自连接与排序}\label{06_sql_queries_zh-ux81eaux8fdeux63a5ux4e0eux6392ux5e8f}

来源：P1 第 50--51 页；P2 第 23 页。

找专业相同的不同学生配对：

\Needspace{105pt}
\begin{lstlisting}[language=SQL]
SELECT a.sid AS sid1, b.sid AS sid2
FROM Student AS a
JOIN Student AS b ON a.major = b.major
WHERE a.sid < b.sid
ORDER BY sid1, sid2;
\end{lstlisting}

按本例学号的通常字符顺序，结果为 \texttt{(\allowbreak{}S1,\allowbreak{}S3)\allowbreak{}}、\texttt{(\allowbreak{}S2,\allowbreak{}S4)\allowbreak{}}。\texttt{\textless{}} 排除自己和反向重复；只有 \texttt{\textless{}\textgreater{}} 会保留正反两种顺序。

\Needspace{75pt}
\begin{lstlisting}[language=SQL]
SELECT cid, title, credits
FROM Course
ORDER BY credits DESC, cid ASC;
\end{lstlisting}

结果顺序为 C1、C3、C2。只按 credits 排序时，C1 与 C3 同为 4，二者先后未确定；若题目要求稳定先后，要增加明确的次级排序键。排序不等于去重，也不改变数据在逻辑关系中的身份。

\subsection{7. 聚合：COUNT、SUM、AVG 与 NULL}\label{06_sql_queries_zh-ux805aux5408countsumavg-ux4e0e-null}

来源：P2 第 12--16 页；补充核对 \href{https://www.postgresql.org/docs/18/functions-aggregate.html}{PostgreSQL 聚合文档}。

\Needspace{120pt}
\begin{lstlisting}[language=SQL]
SELECT COUNT(*) AS n_rows,
       COUNT(score) AS n_scored,
       COUNT(DISTINCT score) AS n_distinct_scores,
       SUM(score) AS total_score,
       AVG(score) AS mean_score
FROM Enrolment;
\end{lstlisting}

本例为 \texttt{(\allowbreak{}6,\allowbreak{}3,\allowbreak{}2,\allowbreak{}250,\allowbreak{}250/\allowbreak{}3)\allowbreak{}}，平均值约 83.33，显示精度取决于数值类型。

\begin{longtable}[]{@{}ll@{}}
\toprule\noalign{}
写法 & 数什么或如何处理 NULL？ \\
\midrule\noalign{}
\endhead
\bottomrule\noalign{}
\endlastfoot
\texttt{COUNT(\allowbreak{}*)\allowbreak{}} & 输入行数，包括 score 缺失的行 \\
\texttt{COUNT(\allowbreak{}score)\allowbreak{}} & score 非空的行数，保留重复值的次数 \\
\texttt{COUNT(\allowbreak{}DISTINCT\ score)\allowbreak{}} & 不同非空 score 值的个数 \\
\texttt{SUM/\allowbreak{}AVG/\allowbreak{}MIN/\allowbreak{}MAX(\allowbreak{}score)\allowbreak{}} & 基于非空输入值计算 \\
\end{longtable}

AVG 的分母是本例 3 个非空成绩，不是 6 次选课。\texttt{AVG(\allowbreak{}DISTINCT\ score)\allowbreak{}} 则为 \texttt{(\allowbreak{}80+90)\allowbreak{}/\allowbreak{}2=85}，已经改变统计问题，不是一般的去重修复。

\subsubsection{空输入与全 NULL}\label{06_sql_queries_zh-ux7a7aux8f93ux5165ux4e0eux5168-null}

无 GROUP BY 的汇总在没有匹配行时仍产生一行：COUNT 为 0，SUM、AVG、MIN、MAX 为 NULL。若匹配行存在但 score 全为 NULL，COUNT(*) 仍数这些行，COUNT(score)=0，其他四种 score 聚合仍为 NULL。

例如筛选 \texttt{sid=\textquotesingle{}S3\textquotesingle{}} 后，COUNT(*)=1、COUNT(score)=0、AVG(score)=NULL。NULL 不是自动补零；是否把缺失汇总显示成 0，需要业务含义支持。

\subsection{8. GROUP BY 与 HAVING：每组一行}\label{06_sql_queries_zh-group-by-ux4e0e-havingux6bcfux7ec4ux4e00ux884c}

来源：P2 第 17--22 页。

\Needspace{120pt}
\begin{lstlisting}[language=SQL]
SELECT e.sid, COUNT(*) AS n_enrolments,
       COUNT(e.score) AS n_scored, AVG(e.score) AS mean_score
FROM Enrolment AS e
GROUP BY e.sid
HAVING COUNT(*) >= 2
ORDER BY e.sid;
\end{lstlisting}

\begin{longtable}[]{@{}llll@{}}
\toprule\noalign{}
sid & n\_enrolments & n\_scored & mean\_score \\
\midrule\noalign{}
\endhead
\bottomrule\noalign{}
\endlastfoot
S1 & 3 & 2 & 85 \\
S2 & 2 & 1 & 80 \\
\end{longtable}

WHERE 筛选输入行，HAVING 筛选已经形成的组。若先加 \texttt{WHERE\ score\ IS\ NOT\ NULL}，再 COUNT(*)，数的就是有成绩的记录，S2 不再满足至少两行，不能声称与原查询相同。

按本课可直接采用的书写规则：SELECT 中非聚合的列都明确列入 GROUP BY；不要查询 \texttt{sid,\allowbreak{}cid,\allowbreak{}AVG(\allowbreak{}score)\allowbreak{}} 却只按 sid 分组，从一个学生的多门课中任意挑 cid。某些实现存在基于函数依赖的例外，不影响先按这条规则写清楚的做法。

\subsubsection{统计对象与零记录}\label{06_sql_queries_zh-ux7edfux8ba1ux5bf9ux8c61ux4e0eux96f6ux8bb0ux5f55}

本例 S1 有 3 次选课，却只有 2 门不同课程；数课程可用 \texttt{COUNT(\allowbreak{}DISTINCT\ cid)\allowbreak{}}。连接 Enrolment 后统计课程数量时，同一课程也可能被多名学生重复计入。必须先回答``数选课记录，还是数不同课程''。

本查询不会为无选课的 S4 生成零计数组。若需要所有学生含零计数，须另设计覆盖全部学生的查询，通常涉及外连接；不能用当前内连接/分组结果冒充已经包含 S4。

\subsection{9. 子查询：先区分标量与多行结果}\label{06_sql_queries_zh-ux5b50ux67e5ux8be2ux5148ux533aux5206ux6807ux91cfux4e0eux591aux884cux7ed3ux679c}

来源：P2 第 2--5 页；基数边界补充见 \href{https://www.postgresql.org/docs/18/functions-subquery.html}{PostgreSQL 子查询文档}。

\subsubsection{标量子查询}\label{06_sql_queries_zh-ux6807ux91cfux5b50ux67e5ux8be2}

\Needspace{75pt}
\begin{lstlisting}[language=SQL]
SELECT cid, title
FROM Course
WHERE credits > (SELECT AVG(credits) FROM Course);
\end{lstlisting}

总体平均学分是 11/3，结果包含 C1、C3。这里的内层聚合保证一列一行，适合当一个值使用。

标量位置要求一列且至多一行：零行相当于 NULL，多于一行在 PostgreSQL/标准标量语义下是错误，不是自动取第一行。例如把内层替换成 \texttt{SELECT\ credits\ FROM\ Course\ WHERE\ dept=\textquotesingle{}CS\textquotesingle{}} 会返回两行，不能直接当标量比较。这个 AVG 子查询不需要引用外层别名。

\subsubsection{IN：测试一个值是否在结果中}\label{06_sql_queries_zh-inux6d4bux8bd5ux4e00ux4e2aux503cux662fux5426ux5728ux7ed3ux679cux4e2d}

\Needspace{135pt}
\begin{lstlisting}[language=SQL]
SELECT s.sid, s.sname
FROM Student AS s
WHERE s.sid IN (
    SELECT e.sid
    FROM Enrolment AS e
    WHERE e.score >= 85
);
\end{lstlisting}

本例只有 S1。单值 IN 的子查询应返回一列；右侧重复值不会让同一外层行重复输出。IN 的逻辑与 NULL 的关系在第 11 节说明。

\subsection{10. EXISTS、NOT EXISTS 与相关子查询}\label{06_sql_queries_zh-existsnot-exists-ux4e0eux76f8ux5173ux5b50ux67e5ux8be2}

来源：P2 第 6 页；Tutorial 2 主题列有 NOT EXISTS。

相关子查询引用外层当前行，例如内层的 \texttt{e.\allowbreak{}sid=s.\allowbreak{}sid}。可按``对每个外层学生问一个存在性问题''理解，不代表数据库一定逐行执行完整子查询。

\Needspace{135pt}
\begin{lstlisting}[language=SQL]
SELECT s.sid, s.sname
FROM Student AS s
WHERE EXISTS (
    SELECT 1
    FROM Enrolment AS e
    WHERE e.sid = s.sid AND e.score IS NULL
);
\end{lstlisting}

本例返回 S1、S2、S3。EXISTS 关心子查询是否有行，不关心 SELECT 列表中的 1；这个普通存在性查询即使有多个匹配，也只为每个外层学生返回一次。外层自己若已有重复行，EXISTS 不会额外替它去重。

\Needspace{135pt}
\begin{lstlisting}[language=SQL]
SELECT s.sid, s.sname
FROM Student AS s
WHERE NOT EXISTS (
    SELECT 1
    FROM Enrolment AS e
    WHERE e.sid = s.sid
);
\end{lstlisting}

本例只有 S4。``从未''要检查\textbf{没有任何匹配行}，不是找一行 \texttt{cid\ \textless{}\textgreater{}\ \textquotesingle{}C1\textquotesingle{}}。若忘记相关条件，内层只要全表存在任意记录，就会影响所有外层学生。

\subsection{11. NOT IN 遇到 NULL：反例要会解释}\label{06_sql_queries_zh-not-in-ux9047ux5230-nullux53cdux4f8bux8981ux4f1aux89e3ux91ca}

来源：教学计划第 2 页 Tutorial 2；补充依据 \href{https://www.postgresql.org/docs/18/functions-subquery.html\#FUNCTIONS-SUBQUERY-NOTIN}{PostgreSQL IN/NOT IN 说明}。

对非空数值，\texttt{80\ NOT\ IN\ (\allowbreak{}90,\allowbreak{}NULL)\allowbreak{}} 等价于 \texttt{80\textless{}\textgreater{}90\ AND\ 80\textless{}\textgreater{}NULL}，结果为 UNKNOWN，WHERE 不保留它；\texttt{90\ NOT\ IN\ (\allowbreak{}90,\allowbreak{}NULL)\allowbreak{}} 则因已找到相等值而为 FALSE。\textbf{不是任何含 NULL 的 NOT IN 都产生同一个结果。}

\begin{longtable}[]{@{}
  >{\raggedright\arraybackslash}p{(\linewidth - 2\tabcolsep) * \real{0.5000}}
  >{\raggedright\arraybackslash}p{(\linewidth - 2\tabcolsep) * \real{0.5000}}@{}}
\toprule\noalign{}
\begin{minipage}[b]{\linewidth}\raggedright
情况
\end{minipage} & \begin{minipage}[b]{\linewidth}\raggedright
单值比较的结果
\end{minipage} \\
\midrule\noalign{}
\endhead
\bottomrule\noalign{}
\endlastfoot
右侧存在相等值 & IN 为 TRUE，NOT IN 为 FALSE \\
没有相等值，但比较涉及 NULL，且右侧非空 & 可能为 UNKNOWN，不能当作确定``不在'' \\
右侧为空集 & IN 为 FALSE，NOT IN 为 TRUE \\
\end{longtable}

对本例这个``找无选课学生''的特定问题，下面写法安全：

\Needspace{75pt}
\begin{lstlisting}[language=SQL]
SELECT s.sid, s.sname
FROM Student AS s
WHERE s.sid NOT IN (SELECT e.sid FROM Enrolment AS e);
\end{lstlisting}

原因是两边 sid 都保证非空；本例仍返回 S4。若换成可空列，这个推理就不成立。删除右侧 NULL 也不能自动解决左侧 NULL 的含义；使用相关 NOT EXISTS 时，同样要明确 NULL 键应如何匹配，不能声称它与 NOT IN 在任意空值条件下都等价。

\subsection{12. UNION、INTERSECT、EXCEPT 与身份粒度}\label{06_sql_queries_zh-unionintersectexcept-ux4e0eux8eabux4efdux7c92ux5ea6}

来源：P2 第 9 页；与 05 的集合运算对应。

\Needspace{165pt}
\begin{lstlisting}[language=SQL]
SELECT e.sid
FROM Enrolment AS e
JOIN Course AS c ON c.cid = e.cid
WHERE c.dept = 'CS'
INTERSECT
SELECT e.sid
FROM Enrolment AS e
JOIN Course AS c ON c.cid = e.cid
WHERE c.dept = 'Math';
\end{lstlisting}

结果为 S2。两侧课程条件不要求同一条选课记录同时属于两个院系；它们先分别求学生集合，再按身份求交。

\begin{longtable}[]{@{}ll@{}}
\toprule\noalign{}
运算 & 默认行为 \\
\midrule\noalign{}
\endhead
\bottomrule\noalign{}
\endlastfoot
UNION & 合并并去除输出重复 \\
UNION ALL & 合并并保留重复次数 \\
INTERSECT & 共有元组，去重 \\
EXCEPT & 左边有、右边没有的元组，去重，方向重要 \\
\end{longtable}

输入列数必须一致，对应位置类型要相容；SQL 集合运算按列位置对应，不要求像本课命名关系代数那样两边列名完全相同。INTERSECT ALL、EXCEPT ALL 另有多重集次数语义，基础复习先与默认去重版本区别，不把 ALL 理解成无作用的装饰。

若先在两边选 sname 再求交，会让 S1 的 CS 记录与 S3 的 Math 记录因同名 Amy 被错误合并；应先求 sid，再连接或用 IN/EXISTS 取得姓名。

\subsection{13. ``每个目标都满足''：没有缺失目标}\label{06_sql_queries_zh-ux6bcfux4e2aux76eeux6807ux90fdux6ee1ux8db3ux6ca1ux6709ux7f3aux5931ux76eeux6807}

来源：P2 第 10--11 页；下例为将 05 方法写成 SQL 的自编解释，不是 Tutorial 逐题答案。

找选过每门 CS 课程的学生，可以从 Student 候选出发，排除存在某门 CS 课却找不到该学生选课记录的情况：

\Needspace{210pt}
\begin{lstlisting}[language=SQL]
SELECT s.sid, s.sname
FROM Student AS s
WHERE NOT EXISTS (
    SELECT 1
    FROM Course AS c
    WHERE c.dept = 'CS'
      AND NOT EXISTS (
          SELECT 1
          FROM Enrolment AS e
          WHERE e.sid = s.sid AND e.cid = c.cid
      )
);
\end{lstlisting}

本例只有 S1。内层寻找某个学生--课程配对；中层寻找缺失配对的目标课程；外层要求这样的目标不存在。跨学期选过即可，是本例明确采用的含义；若要求某个 term，须在目标与选课条件中明确对应的范围。

这要求覆盖所有目标，允许额外选非 CS 课程；重复选同一门课也不能代替缺失的另一门课。目标来自完整 Course，不能仅从已有人选过的课程产生。

若假设数据库中没有 CS 课程，此式返回所有 Student，包括 S4，因为不存在可缺失的目标。若题意要求至少有一个目标，要另加目标非空要求。课件 P2 第 11 页从 Sells 中枚举候选，空目标时只返回在 Sells 中出现过的 bar；这与本例从完整 Student 枚举候选的范围不同，应联系 05 的空除数讨论。

\subsection{14. 课件其余内容与本轮范围}\label{06_sql_queries_zh-ux8bfeux4ef6ux5176ux4f59ux5185ux5bb9ux4e0eux672cux8f6eux8303ux56f4}

\subsubsection{ANY / ALL【补充辨识】}\label{06_sql_queries_zh-any-allux8865ux5145ux8fa8ux8bc6}

来源：P2 第 7--8 页。教学计划 Week 6 列高级子查询及集合比较，本节只帮助识别记号，不据此新增 Test 1 深度要求。

\texttt{x\ \textgreater{}\ ANY\ (subquery)} 问是否至少一个比较为真；\texttt{x\ \textgreater{}\ ALL\ (subquery)} 问是否所有比较都为真。空右侧时 ANY 为 FALSE、ALL 为 TRUE。含 NULL 时还可能 UNKNOWN；不能无条件将 ALL 改写成与 MAX 比较，因为空集或 NULL 可使语义不同。

\subsubsection{DDL、数据修改、视图与 PL/pgSQL}\label{06_sql_queries_zh-ddlux6570ux636eux4feeux6539ux89c6ux56feux4e0e-plpgsql}

P2 第 25--52 页还讲 CREATE TABLE、主外键、INSERT、DELETE、UPDATE、视图和结构修改。这些可联系 03--04 理解，但本章不据文件位置将它们全部算作 Test 1 查询重点：

\begin{itemize}
\tightlist
\item
  PRIMARY KEY 结合唯一性与非空；FOREIGN KEY 并不自动要求该列非空，1:1 也不能仅靠一个普通外键保证。
\item
  DELETE 删除记录，DROP TABLE 删除表结构；UPDATE 没有 WHERE 会作用于全表。这里只解释命令，不执行任何修改。
\item
  普通视图是查询定义的接口，不应将课件``does not take up space''绝对理解为连定义元数据都不占空间，也不等于物化视图保存结果。
\item
  PL/pgSQL、索引、事务、规范化等超出本篇，不展开。
\end{itemize}

\subsection{15. 写查询的检查顺序与复习清单}\label{06_sql_queries_zh-ux5199ux67e5ux8be2ux7684ux68c0ux67e5ux987aux5e8fux4e0eux590dux4e60ux6e05ux5355}

\begin{enumerate}
\def\labelenumi{\arabic{enumi}.}
\tightlist
\item
  \textbf{确定输出粒度}：每个学生、每次选课、不同课程，还是每组一行？
\item
  \textbf{先核对模式和完整连接键}：尤其 sid、cid、term；区分学生专业与课程院系。
\item
  \textbf{写行条件并处理 NULL}：WHERE 只留 TRUE；选择 IS NULL、IN、EXISTS 等合适形式。
\item
  \textbf{明确重复值要求}：需要 DISTINCT 吗？DISTINCT 是在掩盖错误连接，还是符合输出粒度？
\item
  \textbf{分组后再筛组}：SELECT 中的普通列明确分组；COUNT(*)、COUNT(col)、COUNT(DISTINCT col) 不混用。
\item
  \textbf{检查子查询形状与相关条件}：标量是否至多一行？存在性是否真的关联当前外层对象？
\item
  \textbf{检查边界}：同名、重修、未出分、无匹配、重复匹配、空目标；需要顺序时加 ORDER BY。
\end{enumerate}

\begin{longtable}[]{@{}ll@{}}
\toprule\noalign{}
易混点 & 正确区分 \\
\midrule\noalign{}
\endhead
\bottomrule\noalign{}
\endlastfoot
UNKNOWN / FALSE & 逻辑不同；WHERE 都不保留 \\
DISTINCT / GROUP BY & 输出元组去重 / 形成统计组，不可机械互换 \\
WHERE / HAVING & 组形成前筛行 / 聚合后筛组 \\
COUNT(*) / COUNT(score) & 所有输入行 / 非空 score 行 \\
EXISTS / 标量子查询 & 是否有行 / 把结果作为一个值 \\
NOT IN / NOT EXISTS & 比较值集合 / 检查相关匹配；NULL 下不无条件等价 \\
``所有'' / ``只有'' & 覆盖每个目标 / 排除目标之外的联系 \\
SQL 集合运算 / 普通 SELECT & 默认去重 / 默认保留重复 \\
\end{longtable}

以下是复习目标，不是正式模拟卷或官方分值：

\begin{itemize}
\tightlist
\item[$\square$]
  不看笔记写出常用子句的书写顺序与逻辑阶段。
\item[$\square$]
  手算基本查询、DISTINCT、自连接及带次级键的排序。
\item[$\square$]
  用三值逻辑解释为什么未出分记录不满足 \texttt{score\textless{}\textgreater{}90}。
\item[$\square$]
  检查普通连接和跨学期连接的完整条件。
\item[$\square$]
  手算三种 COUNT、全 NULL 与空输入汇总。
\item[$\square$]
  解释 WHERE 与 HAVING 改变统计结果的方式。
\item[$\square$]
  区分标量、IN、EXISTS 子查询的结果要求。
\item[$\square$]
  解释 NOT IN 含 NULL 的反例与非空 sid 的安全情形。
\item[$\square$]
  用学生身份作集合运算，并解释不能过早投影姓名。
\item[$\square$]
  读懂双重 NOT EXISTS 的三层含义及空目标行为。
\end{itemize}

可用于英文简答的表达：

\begin{quote}
WHERE retains only rows for which its condition evaluates to TRUE.
\end{quote}

\begin{quote}
COUNT(*) counts rows, whereas COUNT(column) counts non-NULL values in that column.
\end{quote}

\begin{quote}
WHERE filters input rows; HAVING filters groups after aggregation.
\end{quote}

\begin{quote}
EXISTS tests whether a subquery returns any rows, not the values in its SELECT list.
\end{quote}

\begin{quote}
A NOT IN comparison can evaluate to UNKNOWN when NULL participates and no equal value is found.
\end{quote}

\begin{quote}
A universal requirement can be expressed as the absence of a required target without a matching association.
\end{quote}

\subsection{来源与整理记录}\label{06_sql_queries_zh-ux6765ux6e90ux4e0eux6574ux7406ux8bb0ux5f55}

\begin{itemize}
\tightlist
\item
  教师课件：\href{../../../slides/SQLpart1.pptx}{SQLpart1.pptx}、\href{../../../slides/SQLpart2.pptx}{SQLpart2.pptx}。本篇核心来源为 P1 第 3--58 页、P2 第 2--24 页；已读取文字并重点查看 P1 第 18、31、35--36、51、55、58 页，P2 第 2、6、9、11、15--16、20、22、24 页原图。P2 第 25--52 页仅作范围辨识。
\item
  教学范围：\href{../../../syllabus/CSC5041_Teaching_Plan.pdf}{CSC5041\_Teaching\_Plan.pdf}，第 1 页 Weeks 5--6、第 2 页 Tutorial 2；参考题面：\href{../../../slides/Tutorial\%202.pdf}{Tutorial 2.pdf}，第 1 页。题面 Week 6 与计划 Week 5 标签差异已在 05 记录，本篇不推断改期。
\item
  补充核对：PostgreSQL 18 官方文档的 \href{https://www.postgresql.org/docs/18/functions-comparison.html}{比较}、\href{https://www.postgresql.org/docs/18/functions-subquery.html}{子查询}、\href{https://www.postgresql.org/docs/18/functions-aggregate.html}{聚合}、\href{https://www.postgresql.org/docs/18/queries-table-expressions.html}{表表达式}，于本轮在线读取；用于 NULL、子查询基数及聚合等边界，不代表教师规定了 PostgreSQL 18 版本。
\item
  本版为 2026-10-20 的 LaTeX 重排；示例证据状态沿用原说明，未新增数据库实测。
\item
  输入 SHA-256：SQLpart1 \texttt{\seqsplit{aa28c989fb3674e780418c78c483bba20ff5a50b7da067b0341e60410a68c8ed}}；SQLpart2 \texttt{\seqsplit{acd25cbd4bbd246f8901ada4c437c3f2d00873385961475cd4c941e9089bf45c}}；Tutorial 2 \texttt{\seqsplit{72c016526c1c0ff02c10fcb5f7f964c1184999c684228671025800936a6cabac}}；syllabus 教学计划 \texttt{\seqsplit{31380a7fe6504e46541b6db0bb2f24859c69188bfeb9681dca73354ebeb68295}}。
\end{itemize}
